\section*{Chapitre \ref{ch:g9:arith} --- Arithmétique~: diviseurs et nombres premiers}

\begin{solution}{exo:g9:arith:1}
\omterm{def:g9:arith:divisor}{Diviseurs} de $36$~: $1, 2, 3, 4, 6, 9, 12, 18, 36$.
\omterm{def:g9:arith:divisor}{Diviseurs} de $45$~: $1, 3, 5, 9, 15, 45$.
\omterm{def:g9:arith:divisor}{Diviseurs} de $17$~: $1$ et $17$ seulement ($17$ est premier).
\end{solution}

\begin{solution}{exo:g9:arith:2}
$2\,346$ se termine par $6$~: \omterm{def:g6:wholes:divisible}{divisible} par $2$, pas par $5$. \omterm{def:g1:addition:def}{Somme} des
chiffres $2 + 3 + 4 + 6 = 15$~: \omterm{def:g6:wholes:divisible}{divisible} par $3$, pas par $9$. Les
deux derniers chiffres $46$, et $46 = 4 \times 11 + 2$ n'est pas
\omterm{def:g6:wholes:divisible}{divisible} par $4$~: $2\,346$ n'est pas \omterm{def:g6:wholes:divisible}{divisible} par $4$.
\end{solution}

\begin{solution}{exo:g9:arith:3}
$72 = 2^3 \times 3^2$~; \quad
$150 = 2 \times 3 \times 5^2$~; \quad
$210 = 2 \times 3 \times 5 \times 7$~; \quad
$121 = 11^2$.
\end{solution}

\begin{solution}{exo:g9:arith:4}
$101$~: tester les premiers $p$ avec $p^2 \leq 101$, c'est-à-dire $2,
3, 5, 7$. Aucun ne \omterm{def:g9:arith:divisor}{divise} $101$ (\omterm{def:g3:numbers:evenodd}{impair}, \omterm{def:g1:addition:def}{somme} des chiffres $2$, ne se
termine pas par $0/5$, $101 = 7 \times 14 + 3$)~: $101$ est premier.

$91 = 7 \times 13$~: pas premier.

$143 = 11 \times 13$~: pas premier.
\end{solution}

\begin{solution}{exo:g9:arith:5}
\omterm{def:g9:arith:divisor}{Diviseurs} communs de $48$ et $60$~: \omterm{def:g9:arith:divisor}{diviseurs} de $48$~: $1, 2, 3, 4, 6,
8, 12, 16, 24, 48$~; \omterm{def:g9:arith:divisor}{diviseurs} de $60$~: $1, 2, 3, 4, 5, 6, 10, 12, 15,
20, 30, 60$~; les communs sont $1, 2, 3, 4, 6, 12$, donc
$\gcd(48,60) = 12$.

Par décomposition~: $48 = 2^4 \times 3$ et $60 = 2^2 \times 3 \times 5$~;
premiers communs avec plus petits \omterm{def:g8:powers:def}{exposants}~: $2^2 \times 3 = 12$.
\end{solution}

\begin{solution}{exo:g9:arith:6}
$\gcd(255, 154)$~:
\begin{align*}
255 &= 154 \times 1 + 101, \\
154 &= 101 \times 1 + 53, \\
101 &= 53 \times 1 + 48, \\
53 &= 48 \times 1 + 5, \\
48 &= 5 \times 9 + 3, \\
5 &= 3 \times 1 + 2, \\
3 &= 2 \times 1 + 1, \\
2 &= 1 \times 2 + 0 .
\end{align*}
Dernier \omterm{def:g3:division:remainder}{reste} non nul~: $\gcd(255, 154) = 1$ (ils sont \omterm{def:g9:arith:gcd}{premiers entre
eux}).

$\gcd(1053, 325)$~:
\begin{align*}
1053 &= 325 \times 3 + 78, \\
325 &= 78 \times 4 + 13, \\
78 &= 13 \times 6 + 0 .
\end{align*}
$\gcd(1053, 325) = 13$.
\end{solution}

\begin{solution}{exo:g9:arith:7}
Algorithme d'Euclide~: $588 = 504 \times 1 + 84$~; $504 = 84 \times 6 +
0$~: $\gcd(588, 504) = 84$. Puis
\[
\frac{588}{504} = \frac{588 \div 84}{504 \div 84} = \frac{7}{6},
\]
qui est irréductible.
\end{solution}

\begin{solution}{exo:g9:arith:8}
Le nombre de bouquets doit diviser à la fois $84$ et $126$~; le plus
grand possible est $\gcd(84, 126)$. Décompositions~: $84 = 2^2 \times 3
\times 7$, $126 = 2 \times 3^2 \times 7$, donc le \omterm{def:g9:arith:gcd}{PGCD} est $2 \times 3
\times 7 = 42$. Elle peut faire $42$ bouquets, chacun contenant
$\frac{84}{42} = 2$ roses et $\frac{126}{42} = 3$ tulipes.
\end{solution}

\begin{solution}{exo:g9:arith:9}
On cherche le plus petit \omterm{def:g9:arith:divisor}{multiple} commun. $24 = 2^3 \times 3$ et
$36 = 2^2 \times 3^2$~; en prenant chaque premier avec le \emph{plus
grand} \omterm{def:g8:powers:def}{exposant}~: $\operatorname{ppcm} = 2^3 \times 3^2 = 72$. Les
ferries repartent ensemble $72$ minutes après 8:00, à 9:12.
\end{solution}

\begin{solution}{exo:g9:arith:10}
\emph{1.} Tout \omterm{def:g9:arith:divisor}{diviseur} commun $d$ de $n$ et $n+1$ \omterm{def:g9:arith:divisor}{divise} aussi leur
\omterm{ex:g1:subtraction:difference}{différence} $(n+1) - n = 1$, donc $d = 1$~: $\gcd(n, n+1) = 1$.

\emph{2.} Une \omterm{def:g9:fractions:fraction}{fraction} est irréductible exactement lorsque son
\omterm{def:g4:fractions:def}{numérateur} et son \omterm{def:g4:fractions:def}{dénominateur} sont \omterm{def:g9:arith:gcd}{premiers entre eux}, ce qui est le
cas pour $n$ et $n + 1$ par le point 1.
\end{solution}

\begin{solution}{pb:g9:arith:1}
\textbf{1.} Remplir le $3$ et le verser dans le $5$ (contenus~:
$0/3 \to 3$ dans le grand). Remplir de nouveau le $3$ et le verser
dans le $5$ jusqu'à ce qu'il soit plein~: le grand broc n'accepte
plus que $2$, il reste donc
\[
3 - 2 = 1 \text{ L dans le petit broc.}
\]
Gestes~: remplir le $3$~; verser $3 \to 5$~; remplir le $3$~;
verser $3 \to 5$.

\textbf{2.} Remplir le $5$~; verser dans le $3$ (il reste $2$ dans
le grand)~; vider le $3$~; y verser les $2$~; remplir le $5$~;
verser dans le $3$ jusqu'à le remplir --- il n'accepte que $1$, ce
qui laisse $\mathbf{4}$ L dans le grand broc. Six gestes.

\textbf{3.} Tous~: $1$ (question 1), $2$ (après deux gestes de la
question 2), $3$ et $5$ (simples remplissages), $4$ (question 2),
$6 = 3 + 3$ (un petit broc plein plus $3$ versés dans le grand),
$7 = 5 + 2$, $8 = 5 + 3$ (les deux pleins). Toute quantité entière
de $1$ à $8$ L est mesurable avec le $5$ et le $3$.

\textbf{4.} Au départ, les deux brocs contiennent $0$, un \omterm{def:g9:arith:divisor}{multiple}
de $2$. Remplir fixe un contenu à $6$ ou $4$~: \omterm{def:g3:numbers:evenodd}{pair}. Vider le fixe à
$0$~: \omterm{def:g3:numbers:evenodd}{pair}. Verser déplace de l'eau entre des brocs dont les
contenus étaient \omterm{def:g3:numbers:evenodd}{pairs}, et la quantité versée est une \omterm{ex:g1:subtraction:difference}{différence} de
\omterm{def:g3:numbers:evenodd}{nombres pairs} (la place disponible, ou la quantité présente)~: tous
les contenus restent donc \omterm{def:g3:numbers:evenodd}{pairs} à jamais. Une cible impaire comme
$1$ L est inatteignable.

\textbf{5.} $\gcd(6, 4) = 2$~: seules des quantités paires ---
confirmé par la question 4. $\gcd(5, 3) = 1$~: la loi autorise toute
quantité entière, et la question 3 les a toutes réalisées. Le \omterm{def:g9:arith:gcd}{PGCD}
est exactement l'unité de mesure des brocs.

\textbf{6.} $91 = 1 \times 65 + 26$~; $65 = 2 \times 26 + 13$~;
$26 = 2 \times 13 + 0$~: $\gcd(91, 65) = 13$. Et
$2\,026 = 44 \times 46 + 2$~; $46 = 23 \times 2 + 0$~:
$\gcd(2\,026, 46) = 2$.

\textbf{7.} En versant plusieurs fois le $5$ dans le broc de
$13$~: après deux remplissages, le grand broc contient $10$~; le
troisième remplissage n'y entre qu'à hauteur de $3$, laissant
$5 - 3 = 2$ dans le petit broc --- or le \omterm{def:g3:division:remainder}{reste} de $13$ par $5$
valait $3$, et les quantités $3$ (la place) et $2$ (le reliquat)
sont exactement les nombres d'Euclide ($13 = 2 \times 5 + 3$,
$5 = 1 \times 3 + 2$). En continuant, $3 - 2 = 1$ apparaît~: le
\omterm{def:g3:division:remainder}{reste} suivant de l'algorithme. La fontaine effectue les divisions
d'Euclide avec de l'eau.

\textbf{8.} $\gcd(13, 5) = 1$, donc la loi de la question 5 autorise
toute quantité entière --- et la cascade de la question 7 a
effectivement \omterm{def:g2:mult:def}{produit} $1$ L. Oui.

\textbf{9.} Un \omterm{def:g9:arith:divisor}{diviseur} commun de $n$ et de $2n + 1$ \omterm{def:g9:arith:divisor}{divise}
$2n + 1 - 2 \times n = 1$~: il vaut donc $1$. D'où
$\gcd(n, 2n+1) = 1$ toujours.

\textbf{10.} Bus A~: 7~h~12, 7~h~24, 7~h~36, 7~h~48,
8~h~00\,\dots~; bus B~: 7~h~18, 7~h~36, 7~h~54\,\dots\ Premier
départ commun~: 7~h~36, au bout de $36$ minutes --- le premier
\omterm{def:g9:arith:divisor}{multiple} commun de $12$ et $18$. Loi~:
$36 \times \gcd(12, 18) = 36 \times 6 = 216 = 12 \times 18$. Pour
$5$ et $3$~: premier \omterm{def:g9:arith:divisor}{multiple} commun $15$, et
$15 \times \gcd(5,3) = 15 \times 1 = 15 = 5 \times 3$.

\textbf{11.} Avec un cycle de $17$ ans~: la coïncidence suivante
est le premier \omterm{def:g9:arith:divisor}{multiple} commun de $17$ et de $4$~; comme
$\gcd(17, 4) = 1$, c'est $17 \times 4 = 68$ ans --- les cigales ne
croisent le pic qu'une fois sur quatre sorties. Avec un cycle de
$16$ ans~: $16$ est un \omterm{def:g9:arith:divisor}{multiple} de $4$, donc \emph{chaque} sortie
tombe sur un pic. Une durée de cycle première ne partage aucun
facteur avec un cycle de prédateur plus court, ce qui espace au
maximum les coïncidences~: l'arithmétique comme camouflage.

\textbf{12.} Quatre choix d'\omterm{def:g8:powers:def}{exposant} pour $2$, trois pour $3$, deux
pour $5$~: $4 \times 3 \times 2 = 24$ \omterm{def:g9:arith:divisor}{diviseurs}.

\textbf{13.} Si $m = 2^{a} \times 3^{b} \times \cdots$, alors
$m^2 = 2^{2a} \times 3^{2b} \times \cdots$~: chaque \omterm{def:g8:powers:def}{exposant} est
doublé, donc \omterm{def:g3:numbers:evenodd}{pair}. Or dans $360 = 2^3 \times 3^2 \times 5$, les
\omterm{def:g8:powers:def}{exposants} de $2$ et de $5$ sont \omterm{def:g3:numbers:evenodd}{impairs}~: $360$ n'est pas un carré
parfait.

\textbf{14.} Un \omterm{def:g9:arith:divisor}{diviseur} est son propre partenaire exactement
lorsque $d = \frac nd$, c'est-à-dire $n = d^2$~: seuls les carrés
possèdent un tel \omterm{def:g9:arith:divisor}{diviseur} central. Pour tout autre $n$, les
\omterm{def:g9:arith:divisor}{diviseurs} se répartissent en paires, donc en \omterm{def:g3:numbers:evenodd}{nombre pair}. Ainsi~:
\omterm{def:g3:numbers:evenodd}{nombre impair} de \omterm{def:g9:arith:divisor}{diviseurs} $\Leftrightarrow$ carré parfait.
Vérification~: $36$ a pour \omterm{def:g9:arith:divisor}{diviseurs}
$1, 2, 3, 4, 6, 9, 12, 18, 36$ --- neuf, un \omterm{def:g3:numbers:evenodd}{nombre impair}, et
$36 = 6^2$~; tandis que $360$ en a $24$ (question 12), un \omterm{def:g3:numbers:evenodd}{nombre
pair}, et n'est pas un carré (question 13).

\textbf{15.} Le casier $n$ change d'état une fois par élève $k$
dont le numéro \omterm{def:g9:arith:divisor}{divise} $n$~: en tout, autant de fois que $n$ a de
\omterm{def:g9:arith:divisor}{diviseurs}. Un casier finit \emph{ouvert} lorsque son état a changé
un \omterm{def:g3:numbers:evenodd}{nombre impair} de fois --- d'après la question 14, exactement
lorsque $n$ est un carré parfait. Casiers ouverts~:
$1, 4, 9, 16, 25, 36, 49, 64, 81, 100$ --- il y en a dix.
\end{solution}
